\subsection*{Commutativity and Spectral Decomposition}
The mapping of self-adjoint operators to Kählerian observable functions translates algebraic operator properties into Riemannian and symplectic geometry on the projective Hilbert space. This formalism establishes the geometric spectral theory. We construct observables via geometric projectors, define their spectral properties via critical points, and characterize simultaneous measurability through symplectic involution.

\begin{definition}[Geometric projector]
	\label{def:geometric-projector}
	An observable function $p \in \mathcal{K}(\PH, \R)$ is a \emph{geometric projector} if it is idempotent under the Jordan product: $p \circ_{\hbar} p = p$. By evaluating the $\circ_{\hbar}$ product, this imposes the strict metric constraint
	\begin{equation}
		\frac{\hbar}{2}((p,p)) = p - p^2.
	\end{equation}
\end{definition}

\begin{example}
	For a two-level system $\H=\C^2$, let $\Pi$ be the rank-one orthogonal
	projector, as per Definition \ref{def:orthogonal-projector}, onto a fixed eigenvector, and $p:=\langle\Pi\rangle$. Since
	$\Pi^2=\Pi$, $p$ is a geometric projector in the sense of Definition~\ref{def:geometric-projector}. On the corresponding eigenspace submanifold $\mathbb{P}E_1\subset\PH$, which is a single point in this case, we have $p=1$ and $((p,p))=0$. Instead, on the orthogonal eigenspace $\mathbb{P}E_0$, $p=0$ and again $((p,p))=0$. At every other state $[x]\in\PH$, $p([x])\in(0,1)$ strictly, and the metric constraint $\tfrac{\hbar}{2}((p,p))=p-p^2$ forces $((p,p))>0$ there, in agreement with the remark following Definition~\ref{def:geometric-projector}.
\end{example}

\begin{remark}
	The metric constraint forces $p([x]) \in [0,1]$ for all $[x] \in \PH$, since the Riemann bracket $((p,p)) = g(\Ham{p}, \Ham{p})$ is strictly non-negative. The zero-variance loci where $((p,p)) = 0$ correspond precisely to the regions where the function $p$ attains the absolute values $1$ or $0$. These loci form exactly the totally geodesic eigenspace submanifolds associated with the projector.
\end{remark}

Before stating the spectral decomposition of an observable function, we must geometrically define its spectrum, mapping the algebraic eigenvalues to the critical values of the corresponding Hamiltonian flow.

\begin{definition}[Spectrum of an observable function]
	\label{def:observable-spectrum}
	Let $f = \langle A \rangle \in \mathcal{K}(\PH,\R)$ be an observable function, as in Definition \ref{def:observable-function}, associated with the selfadjoint operator $A \in \BH$. The \emph{spectrum} of $f$, denoted $\sigma(f)$, is the spectrum of the underlying operator $A$, defined in Definition \ref{def:spectrum}. 
\end{definition}

\noindent We can state an equivalent geometrical definition. 

\begin{remark}
	The pure point spectrum of $f$ coincides with the set of critical values of the smooth function $f$ on $\PH$. These spectral values are attained at the states where the Hamiltonian vector field vanishes, \textit{i.e.}, at $\Ham{f} = 0$. 
\end{remark}

With the spectrum and the geometric analogues of orthogonal projectors established, arbitrary finite-dimensional observable functions are constructed via linear combination, translating the algebraic spectral theorem to the geometric domain.

\begin{theorem}[Geometric Spectral Theorem]
	\label{thm:geometric-spectral}
	Let $f \in \mathcal{K}(\PH, \R)$ be an observable function, as in Definition \ref{def:observable-function}, possessing a finite discrete spectrum $\sigma(f) = \{a_1, \dots, a_m\}$. There exists a unique set of geometric projectors $\{p_1, \dots, p_m\} \subset \mathcal{K}(\PH, \R)$ such that
	\begin{enumerate}[label=(\roman*)]
		\item $\sum_{j=1}^m p_j = 1$ (Resolution of the identity).
		\item $p_j \circ_{\hbar} p_k = \delta_{jk} p_j$ (Mutual orthogonality).
		\item $f = \sum_{j=1}^m a_j p_j$ (Spectral decomposition).
	\end{enumerate}
\end{theorem}

The geometric spectral decomposition generalizes to multi-parameter families of observables if and only if their underlying Hamiltonian flows commute. The algebraic-geometric isomorphism of Theorem \ref{thm:lie-jordan-kahler-isomorphism} dictates this condition for jointly measurability.

\begin{definition}[Jointly measurable observables]
	\label{def:jointly-measurable}
	Let $A,B\in\BH$ be selfadjoint operators. The observable functions $f=\langle A\rangle$, $k=\langle B\rangle$, as per Definition~\ref{def:observable-function}, are \emph{jointly measurable} if $[A,B]=0$.
\end{definition}

\begin{remark}
	Definition~\ref{def:jointly-measurable} is well posed. If $\langle A\rangle=\langle A'\rangle$ on $\PH$ for selfadjoint $A,A'\in\BH$, then $A=A'$, by Theorem~\ref{thm:cstar-isomorphism}. Hence $[A,B]=0$ depends only on $f,k$, not on the choice of representing operators.
\end{remark}

To operate entirely within the Kähler manifold formalism, this algebraic requirement must be reformulated using the symplectic structure. The algebraic-geometric isomorphism of Theorem~\ref{thm:lie-jordan-kahler-isomorphism} maps the operator commutator directly to the Poisson bracket, yielding a purely geometric integrability condition.

\begin{proposition}[Geometric commutativity]
	\label{prop:geometric-commutativity}
	Two observable functions $f,k\in\mathcal K(\PH,\R)$ are jointly measurable, as per Definition~\ref{def:jointly-measurable}, if and only if their Poisson bracket vanishes identically, $\poisson{f}{k}=0$.
\end{proposition}

\begin{proof}
	Write $f=\langle A\rangle$, $k=\langle B\rangle$. By
	Proposition~\ref{prop:poisson-comm}, $\poisson{f}{k}=\big\langle\qcomm{A}{B}\big\rangle$, and $\qcomm{A}{B}=\tfrac{1}{i\hbar}[A,B]$ is selfadjoint. By Theorem~\ref{thm:cstar-isomorphism}, the mean value map is injective on selfadjoint operators, so $\big\langle\qcomm{A}{B}\big\rangle\equiv0$ on $\PH$ if and only if $\qcomm{A}{B}=0$, equivalently $[A,B]=0$.
\end{proof}

The evaluation of the Poisson bracket strictly dictates the limits of simultaneous measurability in physical systems, as demonstrated by the non-commuting generators of the angular momentum algebra.

\begin{example}
	Let $f=\langle\sigma_z\rangle$, $k=\langle\Id\rangle\equiv1$. Since
	$\{\sigma_z,\Id\}=0$, Proposition~\ref{prop:poisson-comm} gives
	$\poisson{f}{k}=0$: the two observables commute, and their Hamiltonian
	vector fields, $\Ham{f}$ and $\Ham{k}=0$, trivially commute as vector
	fields. By contrast, for $f=\langle\sigma_x\rangle$, $k=\langle\sigma_y\rangle$,
	$\poisson{f}{k}=\tfrac{2}{\hbar}\langle\sigma_z\rangle\not\equiv0$, so
	the two spin components are not simultaneously measurable, in agreement
	with Proposition~\ref{prop:geometric-commutativity}.
\end{example} 